The following overview is presented as general background knowledge; specific bibliographic citations are omitted in this draft because no retrieved literature excerpts were available.\footnote{No retrieved literature excerpts were available for the Related Work section in this draft. We will add precise citations and tie individual claims to retrieved sources in a subsequent revision.}

Prior work relevant to our problem can be organized into several conceptual strands. First, there is a body of work on long‑form and multi‑paragraph generation that emphasizes planning, hierarchical decomposition, and iterative refinement as means to improve global coherence and manage output length. In practice, such approaches frequently decompose a document into outlines or subunits and apply repeated passes to refine local content and inter‑unit transitions (general background knowledge).

Second, a line of methods concerns structured prompting and document planning. These methods encode higher‑level structure (for example, section and subsection templates, rhetorical role labels, or explicit discourse plans) into prompts or intermediate representations to constrain generation and to reduce downstream editing. Related engineering practices include the use of machine‑readable schemas or templates to enforce output shape and to support automated validation (general background knowledge).

Third, there are approaches that integrate retrieval or external evidence into generation pipelines. Retrieval‑augmented generation (RAG) and related protocols supply candidate snippets or documents to language models at generation time, with subsequent selection or synthesis steps intended to ground assertions in sources. Complementary efforts focus on citation insertion and on post‑hoc verification, aiming to reduce unsupported claims by either constraining the generator to cite supplied evidence or by applying downstream checks that flag assertions lacking provenance (general background knowledge).

Finally, work on citation verification and auditability targets the problem of ensuring that generated factual claims are supported by verifiable sources. Techniques in this area include automated fact‑checking against retrieved passages, provenance tracking for inserted citations, and human‑in‑the‑loop review workflows that exploit structured outputs to make audits efficient (general background knowledge).

Compared to a baseline linear prompting pipeline that produces sections in a fixed sequence, our graph‑aware workflow explicitly represents section dependencies and schedules leaf‑first writing so that local evidence can be gathered and locked in before higher‑level synthesis. We position this design as addressing practical gaps in citation placement, evidence provenance, and post‑generation auditability; the present draft frames these gaps as the targets of our contribution and defers a detailed literature comparison with precise citations to a future revision once retrieval and bibliographic metadata have been incorporated.

In the final paper we will expand this overview into a comprehensive, citation‑anchored related work section that (1) cites concrete papers exemplifying each strand above, (2) contrasts empirical findings and evaluation protocols from prior studies with our experimental setup, and (3) situates our JSON‑constrained, evidence‑aware prompting and assembly pipeline with respect to existing verification and planning techniques.